%%%%%%%%%%%%%%%%%
% CV tailored for SCALIAN - PO & Agents IA
% Positioning: Product Owner IA / Agents IA | Automatisation, gouvernance digitale, delivery & acculturation
%%%%%%%%%%%%%%%%%

\documentclass[8pt,a4paper,ragged2e,withhyper]{altacv}

\geometry{left=1.25cm,right=1.25cm,top=.5cm,bottom=1.cm,columnsep=1.2cm}

\usepackage{paracol}
\usepackage{fontawesome5}
\usepackage{hyperref}

\iftutex 
  \setmainfont{Roboto Slab}
  \setsansfont{Lato}
  \renewcommand{\familydefault}{\sfdefault}
\else
  \usepackage[rm]{roboto}
  \usepackage[defaultsans]{lato}
  \renewcommand{\familydefault}{\sfdefault}
\fi

\definecolor{SlateGrey}{HTML}{2E2E2E}
\definecolor{LightGrey}{HTML}{666666}
\definecolor{DarkPastelRed}{HTML}{8AC0D9}
\definecolor{PastelRed}{HTML}{00B0DA}
\definecolor{GoldenEarth}{HTML}{B4AD0C}

\colorlet{name}{black}
\colorlet{tagline}{PastelRed}
\colorlet{heading}{DarkPastelRed}
\colorlet{headingrule}{GoldenEarth}
\colorlet{subheading}{PastelRed}
\colorlet{accent}{PastelRed}
\colorlet{emphasis}{SlateGrey}
\colorlet{body}{LightGrey}

\definecolor{violet}{HTML}{5A2582}
\definecolor{rouge}{HTML}{F53B3B}
\definecolor{noir}{HTML}{00232C}
\definecolor{bleu}{HTML}{00B0DA}
\definecolor{rose}{HTML}{F831A0}
\definecolor{jaune}{HTML}{FFEC00}
\definecolor{vert}{HTML}{34AD6C}

\renewcommand{\namefont}{\Huge\rmfamily\bfseries}
\renewcommand{\personalinfofont}{\footnotesize}
\renewcommand{\cvsectionfont}{\LARGE\rmfamily\bfseries}
\renewcommand{\cvsubsectionfont}{\large\bfseries}
\renewcommand{\cvItemMarker}{{\small\textbullet}}
\renewcommand{\cvRatingMarker}{\faCircle}

\input{../_shared/pubs-num.tex}

% auto_tex_manager layout tuning
\emergencystretch=3em
\hfuzz=60pt
\hbadness=9999
\sloppy

\begin{document}

\name{BOILEAU Guillaume}
\tagline{Product Owner IA / Agents IA | Automatisation, gouvernance digitale, delivery \& acculturation}

\photoL{2.8cm}{../_shared/moi}

\personalinfo{%
  \email{guillaume.boileaupro@gmail.com}
  \phone{+33 6 59 94 85 84}
  \location{Alpes-Maritimes, FRANCE}
  \linkedin{guillaume-boileau-891b81265}
  \researchgate{Guillaume-Boileau-2}
  \orcid{0000-0002-3576-6968}
}

\makecvheader
\vspace{-0.3cm}

\AtBeginEnvironment{itemize}{\small}
\columnratio{0.64}

\begin{paracol}{2}

\cvsection{Experience}

\cvevent{\textbf{Software Engineer -- Développement logiciel, intégration \& support opérationnel}}{VEV Platform Services France}{2025 -- 2026}{Valbonne, France}

\textbf{Contexte :} Plateforme logicielle industrielle CPMS connectée à des infrastructures de recharge, avec services applicatifs, données opérationnelles, flux d’échanges, contraintes terrain et besoins d’exploitation.

\textbf{Objectif :} Contribuer au développement, à l’intégration et à la fiabilisation de services logiciels opérationnels, en combinant analyse de flux, investigation d’incidents, documentation technique et coordination des sujets avec les équipes.

\begin{itemize}
  \item \textbf{Delivery logiciel :} Contribution à des composants TypeScript, Angular et Node.js intégrés dans une plateforme opérationnelle.
  \item \textbf{Interface métier / technique :} Analyse des comportements applicatifs au regard des contraintes terrain, flux de données et besoins opérationnels.
  \item \textbf{Flux et automatisation :} Investigation d’échanges FTP, interruptions de flux, incohérences de données et comportements inattendus.
  \item \textbf{Support opérationnel :} Diagnostic d’incidents, analyse des causes, formalisation de constats et contribution aux actions de correction.
  \item \textbf{Validation opérationnelle :} Vérification de la cohérence entre données applicatives, comportement des équipements et attendus fonctionnels.
  \item \textbf{Documentation :} Rédaction de constats, analyses d’anomalies et éléments utiles à la résolution par les équipes techniques.
  \item \textbf{Pratiques collaboratives :} Utilisation de Git, Linux, Zenhub et documentation technique pour suivre les sujets et partager l’information.
\end{itemize}

\divider

\cvevent{\textbf{Ingénieur R\&D senior -- Produits data, IA scientifique \& validation}}{CNES}{2023 -- 2025}{Nice, France}

\textbf{Contexte :} Développement logiciel et activités R\&D pour le segment sol de la mission spatiale LISA ESA/NASA, avec chaînes de données mission L0--L3, simulation, intégration de modèles, validation technique et environnement CNES/ESA.

\textbf{Objectif :} Structurer des workflows et outils permettant d’intégrer, comparer, valider et exploiter des modèles et données hétérogènes, avec une logique produit, traçabilité, documentation et revue collaborative.

\begin{itemize}
  \item \textbf{Vision produit technique :} Définition d’outils Python répondant à des besoins d’intégration, de comparaison, de validation et d’exploitation de données mission.
  \item \textbf{Plateforme Python :} Développement de CATpyLISA pour structurer des workflows de données, modèles, analyses et résultats.
  \textcolor{DarkPastelRed}{\href{https://gitlab.oca.eu/gboileau/lisa_l3_tools}{\bf GitLab: CATpyLISA}}
  \item \textbf{Chaînes de données :} Structuration de données mission, configurations, produits intermédiaires, résultats, métriques et niveaux L0--L3.
  \item \textbf{Priorisation technique :} Analyse des besoins, contraintes, écarts, dépendances et points bloquants pour orienter les développements.
  \item \textbf{Validation / IVVQ :} Définition de contrôles de cohérence, règles de comparaison, critères d’acceptation et procédures de vérification.
  \item \textbf{Documentation \& gouvernance :} Formalisation des hypothèses, méthodes, résultats, limites, décisions et éléments de revue.
  \item \textbf{Acculturation technique :} Partage de connaissances, vulgarisation de sujets complexes et accompagnement de contributeurs sur les workflows.
  \item \textbf{Coordination transverse :} Interface avec contributeurs scientifiques, logiciels et institutionnels dans un environnement international.
\end{itemize}

\divider

\cvevent{\textbf{Data Scientist -- Machine learning, données capteurs \& aide à la décision}}{Universiteit Antwerpen, Belgium}{2021 -- 2023}{Antwerp, Belgium}

\textbf{Contexte :} Travaux de data science pour instruments de précision, combinant données capteurs, bruit environnemental, modélisation statistique, machine learning, validation de performance et recommandations techniques.

\textbf{Objectif :} Développer des workflows robustes pour analyser des données complexes, comparer des configurations, évaluer des performances et produire des insights exploitables pour la décision.

\begin{itemize}
  \item \textbf{IA / ML appliqués :} Approches de réduction de bruits non stationnaires avec canaux témoins, machine learning et deep learning.
  \item \textbf{Aide à la décision :} Production de figures de mérite, analyses de sensibilité, contrôles de cohérence et recommandations techniques.
  \item \textbf{Pipelines data :} Développement de workflows Python pour analyse de données, statistiques, performance, validation et reporting.
  \item \textbf{Données capteurs :} Analyse de mesures environnementales, propagation, corrélations spatiales et impacts sur la performance système.
  \item \textbf{Modélisation :} Mise en place de pipelines MCMC pour différentes configurations instrumentales et différents niveaux de bruit.
  \item \textbf{Documentation :} Production de résultats traçables, synthèses techniques et supports pour parties prenantes internationales.
\end{itemize}

\divider

\cvevent{\textbf{Ingénieur R\&D junior -- Simulation, signal \& workflows scientifiques}}{ARTEMIS, Observatoire de la Côte d’Azur}{2018 -- 2021}{Nice, France}

\textbf{Contexte :} Travaux de recherche appliquée liés à l’analyse de signaux faibles, à la simulation numérique, au traitement du signal et à la préparation de missions spatiales.

\textbf{Objectif :} Développer des méthodes logicielles d’analyse et de simulation afin de produire des résultats scientifiques robustes, documentés et exploitables.

\begin{itemize}
  \item \textbf{Simulation numérique :} Développement d’approches de simulation pour environnements complexes et grands jeux de données.
  \item \textbf{Traitement du signal :} Méthodes pour analyser et séparer des signaux faibles et superposés.
  \item \textbf{Python scientifique :} Workflows d’analyse, visualisation, comparaison et validation de résultats.
  \item \textbf{Documentation :} Formalisation des méthodes, hypothèses, résultats et conclusions pour revue collaborative.
  \item \textbf{Communication :} Présentation de résultats complexes à des publics scientifiques et techniques.
\end{itemize}

\switchcolumn

\cvsection{Profile}

\textbf{Product Owner IA / Data \& Automation avec un socle fort en Python, IA appliquée, agents, workflows de données, coordination transverse, documentation, validation et acculturation technique.}

Positionnement pour ce rôle : piloter des produits IA et automatisation de l’idéation au MVP, structurer les besoins, prioriser les backlogs, coordonner les parties prenantes, accompagner l’adoption et industrialiser des pratiques de delivery.

\begin{itemize}
  \item Capacité à faire émerger des cas d’usage IA, les cadrer, les traduire en workflows exploitables et accompagner leur adoption.
  \item Expérience en data, IA appliquée, agents/chatbots/MCP, automatisation, validation de modèles et documentation de pratiques.
  \item Forte capacité de coordination entre profils métiers, techniques, data, logiciel, sponsors et contributeurs.
  \item Expérience en structuration de livrables, roadmap technique, critères d’acceptation, documentation et amélioration continue.
  \item Montée rapide ciblée sur Microsoft Power Platform, Copilot Studio, Dataverse, Power Apps et Power Automate.
\end{itemize}

\cvsection{Languages}

\cvskill{English}{4}
\divider
\cvskill{Spanish}{2}
\divider

\medskip

\cvsection{Education}

\cvevent{PhD in Physics}{Université Côte d’Azur}{2018 -- 2021}{Nice, France}

\divider

\cvevent{Selective Graduate Program in Fundamental Physics (Magistère)}{Université Paris-Saclay}{2016 -- 2018}{Orsay, France}

\divider

\cvevent{MSc in Astronomy and Astrophysics}{Observatoire de Paris / Université Paris-Saclay}{2017 -- 2018}{Paris, France}

\divider

\cvevent{BSc in Fundamental Physics}{Université Paris-Saclay}{2015 -- 2016}{Orsay, France}

\cvsection{Professional Training}

\cvevent{Machine Learning in Python with scikit-learn}{FUN MOOC -- Inria}{2026}{France Université Numérique}

\small{
Predictive modelling, supervised learning, model evaluation, cross-validation, metrics, pipelines and practical machine learning workflows with Python and scikit-learn.
}

\divider

\cvevent{AI, software, web and data training}{OpenClassrooms / Coursera / FUN MOOC}{2020 -- 2026}{}

\small{
Python, SQL, Machine Learning, AI, Linux, JavaScript, TypeScript, Angular, Node.js, Django, REST APIs, C/C++, web development, algorithms, reproducible research and software engineering fundamentals.
}

\divider

\cvevent{Power Platform, Copilot Studio \& Dataverse}{Montée en compétence ciblée}{2026}{}

\small{
Objectif : renforcer la pratique Power Apps, Power Automate, Dataverse, Copilot Studio, Teams, SharePoint et gouvernance Microsoft pour des cas d’usage IA et automatisation.
}

\end{paracol}

\cvsection{Projects}

\cvevent{IA appliquée, agents, chatbots \& MCP}{LLMs / Agentic AI / Automation / Productivité}{2024 -- 2026}{}

Exploration et usage de workflows IA pour revue de code, debugging, documentation, synthèse technique, raisonnement incrémental, logique de chatbot, architectures orientées outils et interactions de type MCP.

\textbf{Rôle :} analyse de cas d’usage, conception de workflows, prompting, assistance au développement, documentation, évaluation de pratiques et veille IA.

\textbf{Compétences mobilisées :} LLMs, agents IA, MCP, chatbots, prompting, Python, architecture de workflows, outillage, acculturation, documentation.

\divider

\cvevent{CATpyLISA / LISA Segment Sol -- Données mission, validation \& delivery scientifique}{Mission spatiale LISA ESA/NASA -- CNES/ESA}{2023 -- 2025}{}

Développement logiciel lié au segment sol de LISA : structuration de données mission, niveaux L0--L3, intégration de modèles, comparaison de résultats, validation technique et revue collaborative de workflows scientifiques.

\textbf{Rôle :} développement Python, structuration de données, cadrage des besoins, validation de modèles, analyse d’écarts, documentation technique, coordination avec contributeurs et préparation de revues.

\textbf{Compétences mobilisées :} Python, data workflows, validation, IVVQ, cycle en V, data quality, documentation, coordination transverse, environnement international.

\divider

\cvevent{VEV -- Logiciel opérationnel, flux de données \& support}{Industrial software / Data flows / Connected infrastructure}{2025 -- 2026}{}

Développement et intégration de services logiciels pour une plateforme CPMS connectée à des infrastructures de recharge, avec analyse de flux FTP, investigation d’incidents, vérification de comportements applicatifs, support opérationnel et documentation technique.

\textbf{Rôle :} développement TypeScript, Angular et Node.js, analyse de flux, investigation d’incidents, validation opérationnelle, formalisation de constats et contribution à la fiabilisation.

\textbf{Compétences mobilisées :} TypeScript, Angular, Node.js, FTP, Linux, Git, analyse d’incidents, intégration logicielle, support opérationnel, documentation technique.

\cvsection{Compétences clés}

\vspace{-.3cm}

\cvsubsection{Product Ownership \& delivery}

{\small
\cvtag{Product Ownership -- ramp-up}
\cvtag{Vision produit}
\cvtag{Backlog prioritization}
\cvtag{Roadmap produit}
\cvtag{MVP}
\cvtag{Delivery}
\cvtag{Critères d’acceptation}
\cvtag{Démos}
\cvtag{SteerCo -- awareness}
\cvtag{Gouvernance produit}
\cvtag{Amélioration continue}
\cvtag{Adoption utilisateurs}
\cvtag{Méthodologies delivery}
}

\divider\smallskip
\vspace{-.3cm}

\cvsubsection{Agents IA, GenAI \& automatisation}

{\small
\cvtag{Agents IA}
\cvtag{LLMs}
\cvtag{MCP}
\cvtag{Chatbots}
\cvtag{Prompt Engineering}
\cvtag{Workflows conversationnels}
\cvtag{Automatisation}
\cvtag{GenAI}
\cvtag{Factory IA -- awareness}
\cvtag{Acculturation IA}
\cvtag{Veille IA}
\cvtag{Bonnes pratiques IA}
\cvtag{Adoption IA}
}

\divider\smallskip
\vspace{-.3cm}

\cvsubsection{Microsoft Power Platform}

{\small
\cvtag{Power Platform -- ramp-up}
\cvtag{Copilot Studio -- ramp-up}
\cvtag{Power Apps -- ramp-up}
\cvtag{Power Automate -- ramp-up}
\cvtag{Dataverse -- ramp-up}
\cvtag{Teams}
\cvtag{SharePoint -- awareness}
\cvtag{Outlook -- usage}
\cvtag{Gouvernance Microsoft -- ramp-up}
\cvtag{Sécurité des solutions -- awareness}
\cvtag{Scalabilité -- awareness}
}

\divider\smallskip
\vspace{-.3cm}

\cvsubsection{Data, décisionnel \& reporting}

{\small
\cvtag{Data Science}
\cvtag{Power BI}
\cvtag{Aide à la décision}
\cvtag{Insights opérationnels}
\cvtag{Consolidation de données}
\cvtag{Historisation -- awareness}
\cvtag{Agrégation de données}
\cvtag{Modèles relationnels}
\cvtag{SQL}
\cvtag{Python}
\cvtag{pandas}
\cvtag{NumPy}
\cvtag{Data quality}
}

\divider\smallskip
\vspace{-.3cm}

\cvsubsection{Coordination, gouvernance \& transformation}

{\small
\cvtag{Coordination transverse}
\cvtag{Communication exécutive}
\cvtag{Sponsors métiers}
\cvtag{Directions métiers}
\cvtag{Conduite du changement}
\cvtag{Gouvernance digitale}
\cvtag{Déploiement international}
\cvtag{Documentation process}
\cvtag{Portails méthodologiques}
\cvtag{Formation}
\cvtag{Vulgarisation}
\cvtag{Anglais professionnel}
}

\divider\smallskip
\vspace{-.3cm}

\cvsubsection{Logiciel, outils \& collaboration}

{\small
\cvtag{Python}
\cvtag{TypeScript}
\cvtag{Angular}
\cvtag{Node.js}
\cvtag{REST API}
\cvtag{Git}
\cvtag{GitLab}
\cvtag{Linux}
\cvtag{Docker}
\cvtag{CI/CD}
\cvtag{Confluence}
\cvtag{Jira}
\cvtag{Zenhub}
\cvtag{Agile}
\cvtag{Kanban}
\cvtag{OpenSearch}
\cvtag{Sentry}
}

\end{document}
